% TOP_PROBLEM: {"id": 30003751, "title": "Regularized Density-to-Potential Maps in Density Functional Theory", "category_id": 16, "category": {"id": 16, "name": "physics", "display_name": "Mathematical Physics", "description": "Problems at the intersection of mathematics and physics.", "slug": "physics", "order_index": 16, "created_at": "2026-07-31T15:26:25.671Z"}, "status": "partially_solved", "statusQualification": "The chosen pure-state or ensemble functional, density space and potential dual space are essential. A regularized density-to-potential map is not silently identified with an unregularized map for every representable density."}
% FIRST_POSTED: 2026-10-11
% ENTRY_KIND: statement_only
% SUBMISSION: {"kind": "statement_only", "category": "top_problems", "problem_id": "30003751", "proposal_type": "established", "source": "Mathematical Methods in Quantum Chemistry (2018), Oberwolfach Reports, DOI 10.4171/owr/2018/13, pp. 631--721. Primary report accessed 2026-10-11; relevant PDF pages 2. https://doi.org/10.4171/owr/2018/13; CDFT Part I (catalogue-listed context; not independently consulted in this preparation). https://arxiv.org/abs/2211.16627; CDFT Part II (catalogue-listed context; not independently consulted in this preparation). https://arxiv.org/abs/2303.01357; Periodic regularized inversion (catalogue-listed context; not independently consulted in this preparation). https://journals.aps.org/prb/abstract/10.1103/PhysRevB.111.205143", "status_evidence": "Canonical source status partially_solved retained. Supplied literature triage dated 2026-08-21: Density to potential mappings and Moreau-Yosida regularization now have rigorous formulations and inversion results in important settings, but general representability and the full mapping problem remain open. Primary-report access checked 2026-10-11; this historical report and inherited triage do not certify current resolution of every component.", "status_checked_on": "2026-10-11", "duplicate_check": "Compared canonical ID, title and normalized target with the complete local UnsolvedMath registry; normalized exact-text matches: none; title matches: none. Related or overlapping canonical records: none identified in this batch. Primary source context and the existing OWR statement audit used for scope and duplicate review; an exhaustive mathematical-equivalence/literature search was not performed.", "attribution": "Original proposer not named in the supplied canonical record; attribution remains to the identified Oberwolfach contribution/report, not its editors as invented proposers. Explanatory statement preparation and source/semantic audit assisted by OpenAI GPT-6 Codex; no human mathematical authorship or independent proof verification is claimed."}
% SOURCE_STATEMENT: "Establish rigorous conditions for existence of the density-to-potential map in density functional theory, and use functional-analytic Moreau–Yosida regularization to address the nondifferentiability of the constrained-search Levy–Lieb functional."
% Preparation and final self-audit: CONTRIBUTING.md prompts C and D; source/audit artifacts retained in the local physics_statements_2026_10_11 batch.
% Imported from: physics_statements.tex; UnsolvedMath ID 30003751
% !TeX program = lualatex
\documentclass[11pt]{article}
\usepackage[margin=25mm]{geometry}
\usepackage{fontspec}
\setmainfont{Latin Modern Roman}
\usepackage{amsmath,amssymb,mathtools}
\usepackage{unicode-math}
\setmathfont{Latin Modern Math}
\usepackage{xurl,hyperref}
\hypersetup{colorlinks=true,urlcolor=blue}
\setcounter{secnumdepth}{0}
\setlength{\parindent}{0pt}
\setlength{\parskip}{5pt}
\setlength{\emergencystretch}{3em}
\newcommand{\sref}[2]{\hyperlink{source:#1:#2}{[S#2]}}
\newcommand{\cataloguescope}[1]{\paragraph{Recorded target.}#1}
\newcommand{\R}{\mathbb R}
\newcommand{\C}{\mathbb C}
\newcommand{\N}{\mathbb N}
\newcommand{\Z}{\mathbb Z}
\newcommand{\Q}{\mathbb Q}
\newcommand{\F}{\mathbb F}
\newcommand{\eps}{\varepsilon}
\begin{document}
% UnsolvedMath ID 30003751; source catalogue code OWR-15998-001
\setcounter{section}{30003750}
\section{Regularized Density-to-Potential Maps in Density Functional Theory}\label{30003751}
{\small UnsolvedMath ID 30003751 · Mathematical Physics}
\subsection{Definitions and mathematical statement}

\paragraph{Shared notation.}$\mathbb N=\{1,2,\ldots\}$; the standard number fields have their usual meanings. All additional conventions are stated locally.


\paragraph{Definitions and conventions.}
For prescribed particle density $\rho$, the Levy--Lieb functional is a constrained infimum of kinetic plus interaction energy over quantum states with that density. A density-to-potential map seeks an external potential realizing the density in the relevant variational problem, modulo an additive constant. Given a Banach or Hilbert density space $X$ and a norm, the Moreau--Yosida envelope has the form $F_\varepsilon(\rho)=\inf_{\sigma\in X}\{F(\sigma)+(2\varepsilon)^{-1}\|\rho-\sigma\|_X^2\}$; differentiability depends on the space and convex formulation.

\paragraph{Target retained from the canonical source catalogue.}
\paragraph{Catalogue wording retained for traceability.}
Establish rigorous conditions for existence of the density-to-potential map in density functional theory, and use functional-analytic Moreau–Yosida regularization to address the nondifferentiability of the constrained-search Levy–Lieb functional.

\paragraph{Source scope and qualification.}
The chosen pure-state or ensemble functional, density space and potential dual space are essential. A regularized density-to-potential map is not silently identified with an unregularized map for every representable density.
\subsection{Short English statement}
Specify when densities determine potentials and use rigorous regularization to handle nondifferentiable density functionals.
\subsection{Sources}
\begin{itemize}
\item[\hypertarget{source:30003751:1}{S1}] ULAM.AI and the UnsolvedMath Contributors, UnsolvedMath, record 30003751 (OWR-15998-001). Catalogue curation under CC BY 4.0; definitions and English exposition added locally.
\url{https://huggingface.co/datasets/ulamai/UnsolvedMath}.
\item[\hypertarget{source:30003751:2}{S2}] Mathematical Methods in Quantum Chemistry (2018), Oberwolfach Reports, DOI 10.4171/owr/2018/13, pp. 631--721. Primary report accessed 2026-10-11; relevant PDF pages 2.
\url{https://doi.org/10.4171/owr/2018/13}.
\item[\hypertarget{source:30003751:3}{S3}] CDFT Part I (catalogue-listed context; not independently consulted in this preparation).
\url{https://arxiv.org/abs/2211.16627}.
\item[\hypertarget{source:30003751:4}{S4}] CDFT Part II (catalogue-listed context; not independently consulted in this preparation).
\url{https://arxiv.org/abs/2303.01357}.
\item[\hypertarget{source:30003751:5}{S5}] Periodic regularized inversion (catalogue-listed context; not independently consulted in this preparation).
\url{https://journals.aps.org/prb/abstract/10.1103/PhysRevB.111.205143}.
\end{itemize}
\label{end:30003751}
\end{document}
